(1) 无delta处满足$\psi''=\kappa^2\psi$。偶态在$|x|<a$取$A\cosh(\kappa x)$，$x>a$取$Be^{-\kappa(x-a)}$；奇态内区取$A\sinh(\kappa x)$，外区类似衰减。每个delta处$\psi'(a^+)-\psi'(a^-)=-2mg\psi(a)/\hbar^2=-2\kappa_0\psi(a)$。结合连续性，偶态满足$\kappa[1+\tanh(\kappa a)]=2\kappa_0$，等价于$\kappa=\kappa_0[1+e^{-2\kappa a}]$；奇态满足$\kappa[1+\coth(\kappa a)]=2\kappa_0$，等价于$\kappa=\kappa_0[1-e^{-2\kappa a}]$。
(2) 对偶态，右边$\kappa_0(1+e^{-2\kappa a})$随$\kappa$从$2\kappa_0$连续降至$\kappa_0$，左边严格上升，因此有唯一正交点。奇态写成$1=\kappa_0(1-e^{-2a\kappa})/\kappa$，右边从$2a\kappa_0$随$\kappa$严格降为0，只有$2a\kappa_0>1$时有非零$\kappa$解；等号仅有$\kappa\to0$阈值态而不可归一。
(3) 若$a\kappa_0\gg1$，令小参数$\varepsilon=e^{-2a\kappa_0}$，用不动点展开得$\kappa_e=\kappa_0(1+\varepsilon+O(a\kappa_0\varepsilon^2))$、$\kappa_o=\kappa_0(1-\varepsilon+O(a\kappa_0\varepsilon^2))$。所以$E_o-E_e=\hbar^2(\kappa_e^2-\kappa_o^2)/(2m)=2\hbar^2\kappa_0^2\varepsilon/m+O(a\kappa_0\varepsilon^2)$。劈裂指数按两阱间隧穿距离$2a$变化。
(4) 在$a\to0$时分布势趋$-2g\delta(x)$，单delta束缚态衰减常数为$2mg/\hbar^2=2\kappa_0$，恰由偶态方程得到$\kappa_e\to2\kappa_0$；奇态存在条件$2a\kappa_0>1$终将失效，其能量上升至连续谱阈值而消失。此极限同时核对了势强度与态数，不能简单把两个独立单阱能量相加。
